\documentclass[10pt,letterpaper]{article}
\usepackage[spanish, mexico]{babel}
\usepackage[latin1]{inputenc}
\usepackage{amsmath}
\usepackage{amsfonts}
\usepackage{amssymb}
\usepackage{moreverb}
\usepackage{graphicx}
\usepackage[export]{adjustbox}
\usepackage{color}
\usepackage{fancyheadings}
\usepackage{listings}
\usepackage{algorithmic}
\usepackage{algorithm}
\floatname{algorithm}{Algoritmo}

%\usepackage{caption}
\usepackage{subcaption}

%-----------------------------------------------------------------------
% Text dimensions 
%-----------------------------------------------------------------------
%% This controls the height of the text
\setlength{\textheight}{22cm}
\setlength{\voffset}{-1.5cm}

%% This controls the width of the text
\setlength{\textwidth}{16cm} % Aproxx. ~16cm
\setlength{\hoffset}{-1.5cm}

%% Separation of the footer from the body text
\setlength{\footskip}{1.5cm} 

%%
%% Becuase I don't need marginal notes, then: 
\setlength{\marginparsep}{0pt}
\setlength{\marginparwidth}{0pt}
%%
%% Header is equal to \textwdith
%% This is used when \marginparsep and \marginparwith are not zero
%\addtolength{\headwidth}{\marginparsep}
%\addtolength{\headwidth}{\marginparwidth}

%% We have to substrac some points on even pages and odd pages
%% in order to match headers and footers on even and odd pages
\addtolength{\evensidemargin}{-56pt}
\addtolength{\oddsidemargin}{-5pt}

%-----------------------------------------------------------------------
% Style of page and headings
%-----------------------------------------------------------------------
\pagestyle{fancy}

\lhead[\fancyplain{}{\bf\thepage}]{\fancyplain{}{\bfseries\rightmark}}
\rhead[\fancyplain{}{\bfseries\leftmark}]{\fancyplain{}{\bf\thepage}}
\lfoot{\sf Depto. de Recursos Naturales, IGF--UNAM}
\rfoot{\sf \copyright 2020, LMCS}
\cfoot{}

%\setlength{\headrulewidth}{0.4pt}
%\setlength{\footrulewidth}{0.4pt}

\newcommand{\mathbi}[1]{\textbf{\em #1}}
\newcommand{\mbf}[1]{\mathbf{#1}}
\newcommand{\mbi}[1]{\textbf{\em #1}}
\newcommand{\mbiu}[1]{{_\textbf{\em #1}}}
\newcommand{\mbu}[1]{{^\textbf{\em #1}}}
\newcommand{\mcons}[1]{\mbox{\textsf{#1}}}
\newcommand{\tensor}[1]{\underline{\underline{#1}}}
\newcommand{\Vector}[1]{\mathbf{#1}}
\renewcommand{\arraystretch}{1.75}

\graphicspath{{../Figuras/}}

\lstset{language=python}
\lstset{commentstyle=\textit}

\title{pyNoxtli}

\author{Luis M. de la Cruz}

\begin{document}

\maketitle

\tableofcontents

 \newpage

\section{Modelo matem�tico: Transferencia de calor}

Ecuaci�n general de transferencia de calor:

\begin{equation*}
c_p \rho \frac{\partial T}{\partial t} +
c_p \rho \frac{\partial}{\partial x_j} \left( u_j T \right) -
\frac{\partial }{\partial x_j} \left( \kappa \frac{\partial T}{\partial x_j}\right) = 
S
\end{equation*}
%
donde se usa la convenci�n de Einstein (�ndices repetidos se suman) y se define lo siguiente:
%
\begin{center}
\begin{tabular}[h!]{cp{10cm}c}
S�mbolo &  & Unidades \\
\hline
\multicolumn{3}{c}{\textbf{Par�metros f�sicos}}\\
$c_p$    & Capacidad calor�fica espec�fica. & [J / Kg $^\text{o}$K]\\
$\rho$    & Densidad. & [Kg / m$^3$]\\
$\kappa$ & Conductividad t�rmica. &  [W / m $^\text{o}$K] \\
$S$ & Ganancia (fuente) o p�rdida (sumidero) de calor & [J/m$^3$ s] \\

$\displaystyle \alpha = \frac{\kappa}{c_p \rho}$ & Difusividad t�rmica. & [m$^2$/s] \\ \\
\hline
\end{tabular}
\end{center}

\begin{center}
\begin{tabular}[h!]{cp{10cm}c}
S�mbolo &  & Unidades \\
\hline
\multicolumn{3}{c}{\textbf{Variables independientes}}\\
$x_j$    & Coordenadas cartesianas de la posici�n: $(x_1, x_2, x_3) \equiv (x, y, z)$. & [m] \\
$t$      & Tiempo. & [s] \\
\hline
\multicolumn{3}{c}{\textbf{Variables dependientes}}\\
$T$    & Temperatura. & [$^\text{o}$K] \\
$u_j$ & Componentes de la velocidad: $(u_1, u_2, u_3) \equiv (u_x, u_y, u_z)$. & [m/s] \\
\hline
\end{tabular}
\end{center}

\newpage
\subsection{Conducci�n de calor: estacionaria}


\begin{equation*}
- \frac{\partial }{\partial x_j} \left( \kappa \frac{\partial T}{\partial x_j}\right) = S
\end{equation*}

\begin{enumerate}

\item Conducci�n de calor estacionaria sin fuentes ni sumideros ($S=0$) :

\begin{equation*}
- \frac{\partial }{\partial x_j} \left( \kappa \frac{\partial T}{\partial x_j}\right) = 0
\end{equation*}

\item Conducci�n de calor estacionaria con $\kappa$ = constante:

\begin{equation*}
- \kappa \frac{\partial^2 T}{\partial x_j \partial x_j} = S
\end{equation*}

\item Conducci�n de calor estacionaria sin fuentes ni sumideros ($S=0$), con $\kappa$ = constante:

\begin{equation*}
- \frac{\partial^2 T}{\partial x_j \partial x_j} = 0
\end{equation*}

\end{enumerate}


\subsubsection{Ejemplo: conducci�n de calor estacionaria sin fuentes ni sumideros con $\kappa$ = cte}

Se considera el problema de conducci�n de calor, sin fuentes, en una barra aislada cuyos extremos se
mantienen a temperatura constante de 100�C y 500�C. Calcular la distribuci�n de temperaturas en la barra

\newpage
\subsection{Conducci�n de calor: NO estacionaria}

\begin{enumerate}
\item Conducci�n de calor no estacionaria:

\begin{equation*}
\frac{\partial T}{\partial t} -
\frac{\alpha}{\kappa} \frac{\partial }{\partial x_j} \left( \kappa \frac{\partial T}{\partial x_j}\right) = \frac{\alpha}{\kappa} S
\end{equation*}

\item Conducci�n de calor no estacionaria sin fuentes ni sumideros ($S=0$) :

\begin{equation*}
\frac{\partial T}{\partial t} - 
\frac{\alpha}{\kappa} \frac{\partial }{\partial x_j} \left( \kappa \frac{\partial T}{\partial x_j}\right) = 0
\end{equation*}

\item Conducci�n de calor no estacionaria con $\kappa$ = constante:

\begin{equation*}
\frac{\partial T}{\partial t} - \alpha \frac{\partial^2 T}{\partial x_j \partial x_j} = \frac{\alpha}{\kappa} S
\end{equation*}
	
\item Conducci�n de calor no estacionaria sin fuentes ni sumideros ($S=0$), con $\kappa$ = constante:

\begin{equation*}
\frac{\partial T}{\partial t} - \alpha \frac{\partial^2 T}{\partial x_j \partial x_j} = 0
\end{equation*}

\end{enumerate}

\newpage
\subsection{Convecci�n de calor: estacionaria}

\begin{enumerate}
\item Convecci�n de calor estacionaria:

\begin{equation*}
\frac{\partial}{\partial x_j} \left( u_j T \right) -
\frac{\alpha}{\kappa} \frac{\partial }{\partial x_j} \left( \kappa \frac{\partial T}{\partial x_j}\right) = 
\frac{\alpha}{\kappa} S
\end{equation*}
	
\item Convecci�n de calor estacionaria sin fuentes ni sumideros ($S=0$):

\begin{equation*}
\frac{\partial}{\partial x_j} \left( u_j T \right) -
\frac{\alpha}{\kappa} \frac{\partial }{\partial x_j} \left( \kappa \frac{\partial T}{\partial x_j}\right) = 0
\end{equation*}

\item Convecci�n de calor estacionaria con $\kappa$ = constante:

\begin{equation*}
\frac{\partial}{\partial x_j} \left( u_j T \right) -
\alpha \frac{\partial^2 T}{\partial x_j \partial x_j} = 
\frac{\alpha}{\kappa} S
\end{equation*}

\item Convecci�n de calor estacionaria sin fuentes ni sumideros ($S=0$), con $\kappa$ = constante:

\begin{equation*}
\frac{\partial}{\partial x_j} \left( u_j T \right) -
\alpha \frac{\partial^2 T}{\partial x_j \partial x_j} = 0
\end{equation*}

\end{enumerate}

\newpage
\subsection{Convecci�n de calor: NO estacionaria}

\begin{enumerate}
\item Convecci�n de calor no estacionaria:

\begin{equation*}
\frac{\partial T}{\partial t} +
\frac{\partial}{\partial x_j} \left( u_j T \right) -
\frac{\alpha}{\kappa} \frac{\partial }{\partial x_j} \left( \kappa \frac{\partial T}{\partial x_j}\right) = 
\frac{\alpha}{\kappa} S
\end{equation*}
	
\item Convecci�n de calor estacionaria sin fuentes ni sumideros ($S=0$):

\begin{equation*}
\frac{\partial T}{\partial t} +
\frac{\partial}{\partial x_j} \left( u_j T \right) -
\frac{\alpha}{\kappa} \frac{\partial }{\partial x_j} \left( \kappa \frac{\partial T}{\partial x_j}\right) = 0
\end{equation*}

\item Convecci�n de calor estacionaria con $\kappa$ = constante:

\begin{equation*}
\frac{\partial T}{\partial t} +
\frac{\partial}{\partial x_j} \left( u_j T \right) -
\alpha \frac{\partial^2 T}{\partial x_j \partial x_j} = 
\frac{\alpha}{\kappa} S
\end{equation*}

\item Convecci�n de calor estacionaria sin fuentes ni sumideros ($S=0$), con $\kappa$ = constante:

\begin{equation*}
\frac{\partial T}{\partial t} +
\frac{\partial}{\partial x_j} \left( u_j T \right) -
\alpha \frac{\partial^2 T}{\partial x_j \partial x_j} = 0
\end{equation*}

\end{enumerate}
\newpage

\end{document}